% !PRINTED 330
\weberpage{330}
%<<CONT
worin $h$, $a$, $b$, $c$ Constanten sind, von denen $h$ von Null verschieden
 ist, und $f_3$ eine binäre cubische Form von $x_1, x_2$.

Wenn wir hiernach die nach absteigenden Potenzen von $x_3$
geordnete Function $H$ berechnen:

\weq{2}{H = x_3^{3n-6}\,H_0 + x_3^{3n-7}\,H_1 + \cdots,}

so findet sich

\weq{3}{\begin{aligned} H_0 &= -\,2\,(n-1)^2\,h^2\,c \\ H_1 &= -(n-1)^2\,h^2\,f_3''\,(x_2,x_2) + 4\,(n-1)(n-2)\,h\,(b^2-a\,c)\,x_1. \end{aligned}}

Die durch die Gleichung $H=0$ dargestellte Curve heisst
die \gesperrt{Hesse'sche Curve} der Curve $f$. Die Formeln (2) und (3)
zeigen, dass die Hesse'sche Curve dann und nur dann durch den
Punkt $\xi_3$ hindurchgeht, wenn $H_0=0$, also $c=0$, d.\,h. nach \S.\,86,
wenn dieser Punkt ein Wendepunkt ist. Aus dem Ausdrucke
für $H_1$ kann man noch weiter schliessen, dass die Wendetangente
$x_1=0$ dann und nur dann zugleich Tangente der Curve $H=0$
ist, wenn $f_3''\,(x_2,x_2)$ und folglich auch $f_3$ selbst durch $x_1$ theilbar
ist. In diesem Falle wäre, wie aus (3) hervorgeht, $x_1=0$ eine
im Punkte $\xi_3$ \gesperrt{vierpunktig berührende Tangente}. Dieser
Fall kann bei den Curven dritter Ordnung nur dann vorkommen,
wenn $f$ durch $x_1$ theilbar ist, also niemals bei einer irreduciblen
Curve dritter Ordnung.

Die \gesperrt{Hesse'sche} Curve geht ausser durch die Wendepunkte
noch durch die singulären Punkte der Curve $f$, was hier nicht
weiter untersucht werden soll.

Wenn wir also unsere Betrachtungen auf Curven $f$ ohne
singulären Punkt beschränken, so können wir sagen, dass jeder
Schnittpunkt der Curve $f$ mit ihrer \gesperrt{Hesse'schen} Curve einen
Wendepunkt der Curve $f$ giebt, und dass diese Curve $f$ keine
anderen Wendepunkte hat. Wenn Punkte mit vierpunktig berührender
 Tangente vorkommen, so berühren sich die Curven $f$
und $H$, und wir können solche Punkte auffassen als durch das
Zusammenfallen von zwei Wendepunkten entstanden. Wenden
wir noch das Theorem von Bezout (Bd.\,I, \S.\,49) an, so erhalten
wir das Ergebniss:

\begin{quote}
\gesperrt{Eine Curve $n^{\mathrm{ter}}$ Ordnung ohne singulären Punkt
hat $3\,n\,(n-2)$ Inflexionspunkte;}
\end{quote}

und speciell:

\begin{quote}
\gesperrt{Eine Curve dritter Ordnung ohne singulären Punkt
hat neun getrennt liegende Inflexionspunkte.}
\end{quote}
